\chapter*{Abstract}
\thispagestyle{plain}

This thesis presents AK\&VEN, the online shop of a family business that has sold Korean-made socks from container
70-E at Dordoi Bazaar in Bishkek, Kyrgyzstan, since 2018. At the bazaar the price is agreed rather than read from a
label, families buy for everyone at once, and shops from the regions buy by the case. The shop brings these habits
online instead of replacing them with a fixed price tag: customers mix any socks and the price per pair falls with the
size of the whole bag, plan one set for the whole family, ask for a better price in a chat, and collect their order at
the stall with a pickup code or have it delivered to four countries.

The central technical question is how a language model can be allowed to talk about money without being trusted with
it. The design separates talking from deciding. The assistant, either a rule-based negotiator or a large language
model, only proposes a discount; a small deterministic component, the \texttt{PolicyValidator}, clamps every proposal
to a limit the owners set for each product, which the model never sees. The check is repeated at checkout and by
the database, and every proposal is recorded beside the discount allowed. A model that is confused, flattered or manipulated therefore cannot change what a customer pays.

The system is a single installable progressive web application in English, Russian and Kyrgyz, built with React and
TypeScript, on top of one Spring Boot API with role-based access for customers, staff, owners and a super admin, and a
PostgreSQL database whose constraints repeat the most important business rules. The thesis documents the requirements,
the interface for each role, the architecture, the data model, the pricing and negotiation logic, security, and more
than 470 automated tests. A recorded run of a live model on sixteen simulated customers shows that neither assistant
ever priced a sock below the owners' limit.
